\chapter{Lezione 1 --- Sistemi in retroazione}

\section{Schema generale del sistema in retroazione}

Consideriamo un sistema lineare tempo-invariante descritto nel dominio di Laplace
da un regolatore \(R(s)\) e da un impianto \(G(s)\), con retroazione unitaria negativa.

Indicando con:
\begin{itemize}
    \item \(y^\star(t)\) il riferimento,
    \item \(e(t)\) l'errore,
    \item \(u(t)\) il segnale di controllo,
    \item \(y(t)\) l'uscita,
\end{itemize}
si ha nel dominio di Laplace
\[
E(s)=Y^\star(s)-Y(s),
\qquad
U(s)=R(s)E(s),
\qquad
Y(s)=G(s)U(s).
\]

La funzione d'anello, o funzione di catena diretta, è
\[
L(s)=R(s)G(s).
\]

La funzione di trasferimento in ciclo chiuso dal riferimento all'uscita è
\[
H(s)=\frac{Y(s)}{Y^\star(s)}=\frac{L(s)}{1+L(s)}=\frac{R(s)G(s)}{1+R(s)G(s)}.
\]

Nel seguito si assume che le eventuali dinamiche non visibili dall'esterno
(non osservabili o non raggiungibili) siano stabili, in modo da poter
ricondurre l'analisi della stabilità del sistema alla stabilità del ciclo chiuso.

\section{Obiettivo del progetto del regolatore}

Dato l'impianto \(G(s)\), il problema del controllo consiste nel progettare \(R(s)\)
in modo che il sistema in retroazione soddisfi alcune specifiche.

Le richieste principali si dividono tipicamente in tre gruppi:
\begin{enumerate}
    \item stabilità;
    \item prestazioni statiche, cioè a regime;
    \item prestazioni dinamiche, cioè nel transitorio.
\end{enumerate}

\section{Stabilità}

Una prima richiesta fondamentale è la stabilità BIBO
(\emph{Bounded Input, Bounded Output}):
a ogni ingresso limitato deve corrispondere un'uscita limitata.

In formule:
\[
y^\star \in L_\infty \quad \Longrightarrow \quad y \in L_\infty.
\]

La stabilità BIBO è la proprietà minima richiesta al sistema in ciclo chiuso.
In molti problemi si richiede anche la stabilità asintotica, cioè che i transitori
si estinguano nel tempo.

\section{Prestazioni statiche}

Le prestazioni statiche riguardano l'errore a regime, cioè la differenza tra
riferimento e uscita quando il transitorio si è esaurito.

Definiamo
\[
e_\infty=\lim_{t\to+\infty}\bigl(y^\star(t)-y(t)\bigr),
\]
quando tale limite esiste.

L'errore a regime dipende dal comportamento di \(L(s)\) alle basse frequenze
e, in particolare, dal numero di poli nell'origine presenti nella funzione d'anello.

\subsection{Tipo del sistema}

Si dice che il sistema è di tipo \(\nu\) se la funzione d'anello \(L(s)\)
ha \(\nu\) poli nell'origine, cioè se si può scrivere localmente per \(s\to 0\) come
\[
L(s)\sim \frac{K}{s^\nu},
\qquad \nu=0,1,2,\dots
\]
con \(K>0\).

Il numero di integratori nella catena diretta è quindi ciò che determina
la capacità del sistema di inseguire segnali di riferimento ``lenti''.

\subsection{Ingressi tipici}

I riferimenti canonici sono:
\[
\text{gradino: } r(t)=1(t),
\qquad
\text{rampa: } r(t)=t\,1(t),
\qquad
\text{parabola: } r(t)=\frac{t^2}{2}\,1(t),
\]
le cui trasformate di Laplace sono rispettivamente
\[
R(s)=\frac{1}{s},
\qquad
R(s)=\frac{1}{s^2},
\qquad
R(s)=\frac{1}{s^3}.
\]

\subsection{Errore a regime per gradino, rampa e parabola}

Per sistemi in retroazione unitaria, le regole classiche sono:

\[
\begin{array}{c|ccc}
 & \nu=0 & \nu=1 & \nu=2 \\ \hline
\text{gradino} & \dfrac{1}{1+K_p} & 0 & 0 \\[1.2ex]
\text{rampa} & +\infty & \dfrac{1}{K_v} & 0 \\[1.2ex]
\text{parabola} & +\infty & +\infty & \dfrac{1}{K_a}
\end{array}
\]
dove:
\[
K_p=\lim_{s\to 0}L(s),\qquad
K_v=\lim_{s\to 0}sL(s),\qquad
K_a=\lim_{s\to 0}s^2L(s).
\]

Questa tabella mostra che:
\begin{itemize}
    \item un alto guadagno a bassa frequenza riduce l'errore statico;
    \item la presenza di integratori è essenziale per annullare l'errore
    su riferimenti di tipo rampa o parabola.
\end{itemize}

\subsection{Interpretazione in frequenza}

Le prestazioni statiche si leggono dal comportamento di \(L(j\omega)\) per \(\omega\to 0\).

Se
\[
|L(j\omega)|\gg 1 \qquad \text{per } \omega\to 0,
\]
allora
\[
H(j\omega)=\frac{L(j\omega)}{1+L(j\omega)}\approx 1,
\]
cioè il sistema in ciclo chiuso riproduce bene il riferimento alle basse frequenze.

Questa è la ragione per cui si desidera che la funzione d'anello abbia
modulo elevato in bassa frequenza.

\section{Risposta armonica}

Se il sistema è lineare e asintoticamente stabile, a un ingresso sinusoidale
\[
u(t)=A\sin(\omega t)
\]
corrisponde, a regime, un'uscita sinusoidale alla stessa pulsazione:
\[
y(t)=A\,|H(j\omega)|\,\sin\bigl(\omega t+\varphi(\omega)\bigr),
\]
dove
\[
\varphi(\omega)=\arg H(j\omega).
\]

In forma complessa:
\[
H(j\omega)\in\mathbb{C},
\qquad
H(j\omega)=|H(j\omega)|\,e^{j\varphi(\omega)}.
\]

Il modulo \(|H(j\omega)|\) dice quanto il sistema amplifica o attenua
una sinusoide a quella frequenza; la fase \(\varphi(\omega)\) dice quanto la ritarda o anticipa.

\section{Prestazioni dinamiche}

Le prestazioni dinamiche riguardano il transitorio, cioè il modo in cui il sistema
raggiunge il regime.

\subsection{Sistema del primo ordine}

Un modello elementare è
\[
H(s)=\frac{1}{1+s\tau}.
\]

La risposta al gradino unitario è
\[
y(t)=1-e^{-t/\tau}.
\]

Il parametro \(\tau\) è la costante di tempo:
più \(\tau\) è piccola, più il sistema è rapido.

\subsection{Sistema del secondo ordine}

Un modello standard del secondo ordine è
\[
H(s)=\frac{\omega_n^2}{s^2+2\zeta\omega_n s+\omega_n^2}
=
\frac{1}{1+2\zeta \dfrac{s}{\omega_n}+\dfrac{s^2}{\omega_n^2}}.
\]

I parametri principali sono:
\begin{itemize}
    \item \(\omega_n\): pulsazione naturale;
    \item \(\zeta\): coefficiente di smorzamento.
\end{itemize}

Se \(0<\zeta<1\), la risposta al gradino è oscillante smorzata.
All'aumentare di \(\zeta\), le oscillazioni si riducono.

\subsection{Margine di fase}

Nel progetto in frequenza è molto importante il margine di fase \(M_\varphi\).

Sia \(\omega_c\) la pulsazione di attraversamento in ampiezza, definita da
\[
|L(j\omega_c)|=1
\qquad
\text{(oppure 0 dB)}.
\]

Il margine di fase è
\[
M_\varphi = \pi + \arg L(j\omega_c)
\]
oppure, in gradi,
\[
M_\varphi = 180^\circ + \arg L(j\omega_c).
\]

Interpretazione:
\begin{itemize}
    \item se il margine di fase è troppo piccolo, il sistema è poco smorzato
    o addirittura instabile;
    \item un margine di fase intorno a \(45^\circ\)-\(70^\circ\) è spesso
    associato a un comportamento soddisfacente.
\end{itemize}

In modo euristico, per sistemi dominati dal secondo ordine, si usa spesso la regola:
\[
\zeta \approx \frac{M_\varphi}{100}
\qquad
\text{(con \(M_\varphi\) espresso in gradi, stima grossolana).}
\]

\section{Bassa e alta frequenza}

Il progetto del controllore deve tenere conto sia del comportamento a bassa frequenza
sia di quello ad alta frequenza.

\subsection{Bassa frequenza}

Alle basse frequenze si vuole:
\[
|L(j\omega)| \gg 1,
\]
in modo da avere
\[
H(j\omega)\approx 1.
\]

Questo garantisce buone prestazioni statiche e un buon inseguimento del riferimento.

\subsection{Alta frequenza}

Alle alte frequenze, se
\[
|L(j\omega)|\ll 1,
\]
allora
\[
H(j\omega)\approx L(j\omega).
\]

In genere è desiderabile attenuare le alte frequenze, perché lì si collocano
rumore di misura, componenti non modellate e dinamiche indesiderate.

\section{Disturbi, sensibilità e sensibilità complementare}

Introduciamo ora due disturbi tipici:
\begin{itemize}
    \item \(m\): disturbo di misura, cioè rumore sul segnale misurato;
    \item \(d\): disturbo di carico, applicato in uscita all'impianto.
\end{itemize}

Assumiamo il modello standard:
\[
e(t)=y^\star(t)-\bigl(y(t)+m(t)\bigr),
\]
e che il disturbo \(d(t)\) si sommi all'uscita dell'impianto.

Nel dominio di Laplace si ottiene
\[
Y(s)=T(s)\,Y^\star(s)-T(s)\,M(s)+S(s)\,D(s),
\]
dove:
\[
T(s)=\frac{R(s)G(s)}{1+R(s)G(s)}
\]
è la \emph{sensibilità complementare}, mentre
\[
S(s)=\frac{1}{1+R(s)G(s)}
\]
è la \emph{sensibilità}.

Le due funzioni soddisfano la relazione fondamentale
\[
S(s)+T(s)=1.
\]

\subsection{Significato fisico}

\begin{itemize}
    \item \(T(s)\) governa l'effetto del riferimento e del rumore di misura sull'uscita.
    \item \(S(s)\) governa l'effetto dei disturbi di carico e l'errore residuo.
\end{itemize}

Di conseguenza:
\begin{itemize}
    \item a bassa frequenza si desidera \(S(j\omega)\approx 0\), quindi \(T(j\omega)\approx 1\),
    per rigettare i disturbi lenti e seguire bene il riferimento;
    \item ad alta frequenza si desidera \(T(j\omega)\approx 0\), quindi \(S(j\omega)\approx 1\),
    per non amplificare il rumore di misura.
\end{itemize}

Questa è la classica tensione di progetto del controllo in frequenza:
non si può rendere contemporaneamente piccolo sia \(S\) sia \(T\) alla stessa frequenza,
perché la loro somma è sempre uguale a \(1\).

\section{Filtri ideali}

Gli andamenti qualitativi richiamati nella lezione sono quelli dei filtri ideali.

\subsection{Passa-basso}

Un filtro passa-basso lascia passare le basse frequenze e attenua le alte:
\[
|H(j\omega)| \approx
\begin{cases}
1, & \omega \ll \omega_c,\\
0, & \omega \gg \omega_c.
\end{cases}
\]

\subsection{Passa-alto}

Un filtro passa-alto attenua le basse frequenze e lascia passare le alte:
\[
|H(j\omega)| \approx
\begin{cases}
0, & \omega \ll \omega_c,\\
1, & \omega \gg \omega_c.
\end{cases}
\]

\subsection{Passa-banda}

Un filtro passa-banda lascia passare un intervallo di frequenze compreso tra
\(\omega_1\) e \(\omega_2\), attenuando quelle molto basse e quelle molto alte.

\section{Sintesi della lezione}

La progettazione del regolatore \(R(s)\) si basa quindi su pochi principi chiave:
\begin{enumerate}
    \item rendere stabile il sistema in ciclo chiuso;
    \item garantire prestazioni statiche adeguate attraverso alto guadagno a bassa frequenza
    e opportuno numero di integratori;
    \item garantire prestazioni dinamiche adeguate controllando banda e margine di fase;
    \item attenuare i disturbi mediante le funzioni di sensibilità
    \[
    S(s)=\frac{1}{1+R(s)G(s)},\qquad
    T(s)=\frac{R(s)G(s)}{1+R(s)G(s)}.
    \]
\end{enumerate}